\documentclass[10pt,a4paper]{amsart}
\usepackage[margin=19mm]{geometry}
\usepackage{amsmath,amssymb,mathtools}
\usepackage[scr=rsfs,cal=euler]{mathalfa}
\usepackage{xfrac}
\usepackage[unicode,hidelinks,bookmarks=false]{hyperref}
\newcommand{\ie}{i.e.\ }
\newcommand{\cf}{cf.\ }
\newcommand{\resp}{respectively\ }
\newcommand{\Subsection}{\subsection}
\providecommand{\nobreakdash}{\nobreak\hbox{-}}
\setlength{\emergencystretch}{2em}
\multlinegap0pt
\usepackage{amssymb}
\usepackage{xcolor}
\newtheorem{lemma}{Lemma}
\title{Anchored Peskin Problem}
\author{Achyuta Telekicherla Kandalam, Daniel Spirn}
\date{Source: arXiv:2604.27219v1 (CC BY 4.0)}
\begin{document}
\maketitle

\noindent\emph{Source and licence.} Anchored Peskin Problem. Version arXiv:2604.27219v1 (CC BY 4.0), distributed by arXiv under the Creative Commons Attribution 4.0 International licence, http://creativecommons.org/licenses/by/4.0/, as recorded in the arXiv licence metadata for this exact version and shown on the abstract page https://arxiv.org/abs/2604.27219v1. Full licence notice: \texttt{licenses/arxiv-2604.27219-CC-BY-4.0.txt}.\par\smallskip

\noindent\textit{Editorial scope.} Counted unit: the article's lemma lem:sinequiv (source0.tex lines 3265--3271). The article's complete original proof of it is delivered with it (source0.tex lines 3273--3368, one \texttt{proof} environment of 96 lines). No mathematics is written, repaired or re-derived: the statement and the proof are the article's own lines, in the article's own notation, and no retained line is substituted or re-flowed.  No further source-local context is needed: every cross-reference of every retained line resolves inside the two delivered spans.  Nothing was omitted from any retained span, so the package carries no omission record.  No retained line contains a citation, so the wrapper carries no bibliography and no bibliography entry.  No retained line contains a figure environment, an image-inclusion command or drawing code, so no image is delivered and the package declares no asset dependency.  Editorial formatting only, in addition: an AMS \texttt{amsart} preamble that restates the article's own declarations and macros used by the retained lines, namely the environments \texttt{lemma} on separate counters, together with the standard packages those lines need.  The retained environments print in this document as Lemma 1 in order of appearance, whereas the article prints the same items with the numbering of its own full document.  The complete native source of the article as distributed in the arXiv e-print for this exact version is bundled unchanged as \texttt{sources/199455/source0.tex}.
\begin{lemma} \label{lem:sinequiv}
For all $0 \leq x,y \leq \pi$ we have $\sin\left( {x+y \over 2}\right) \sim  \sin\left( x \right)+ |x-y| $,  in the sense that 
\[
{1\over C}\left( \sin\left(x\right) + |x-y| \right) \leq \sin \left( {x+y \over 2}\right) \leq C \left( \sin\left( x \right)+ |x-y| \right)
\]
with $C = 8\pi$.
\end{lemma}

\begin{proof}
We begin with elementary trigonometric bounds:
\[
\begin{cases}
{2\over \pi } z \leq \sin(z) \leq z &  \hbox{ for } 0 \leq z \leq {\pi \over 2} \\
1 - {2\over \pi } \left( z - {\pi \over 2} \right) \leq \sin(z) \leq \pi - z & \hbox{ for } {\pi \over 2} \leq z \leq {\pi }  \\
{1\over \sqrt{2}}  \leq \cos(z) \leq 1 &  \hbox{ for } 0 \leq z \leq {\pi \over 4}.
\end{cases}
\]
We proceed by cutting $[0,\pi]^2$ into four regions.

\textbf{Case 1.} We first consider $0\leq x,y \leq {\pi \over 4}$.  
$0 \leq {x + y \over 2} \leq {\pi \over 4}$.
We now consider two separate cases.  First suppose $0 \leq x \leq y \leq {\pi \over 4}$ and set $h = y-x$, then $0 \leq h \leq {\pi \over 4}$.  Then 
$x + y = 2x + h$, or ${x + y \over 2} = x + {h\over 2}$.  Returning to our estimate, note that
\[
\sin \left( {x + y\over 2} \right)
= \sin \left( x \right) \cos \left( {h\over 2} \right) + \cos(x) \sin\left( {h\over 2} \right).
\]
Using our earlier bounds, we find
\begin{align*}
\sin \left( {x + y\over 2} \right)
& \leq  \sin\left( x \right)  + {h \over 2}  \leq  \sin\left( x \right) + h ,
\end{align*}
and 
\begin{align*}
\sin \left( {x + y\over 2} \right) 
& \geq  \sin\left(x \right) \cos\left( {\pi \over 8} \right)+ \cos\left({\pi \over 4} \right) {h \over \pi}   \geq {3 \over 4} \sin\left( x\right) + {1 \over \sqrt{2} \pi}h \\
& \geq {1 \over \sqrt{2}\pi} \left( \sin\left( x\right) + h \right).
\end{align*}

We now consider, $0 \leq y \leq x \leq {\pi \over 4}$.  In this case we can write $x + |x-y| = 2x - y$, and so
\begin{align*}
    x &\leq x+ y \leq 2x\\
     x & \leq x + |x - y| \leq 2x.  
\end{align*}
Therefore, $x+ y  \leq 2 x \leq 2 ( 2 x - y)= 2 ( x + |x -y|)$, 
and $x+ y \geq x \geq {1\over 2} ( 2 x - y) = {1\over 2} ( x + |x -y|)$, 
which implies
\begin{align*}
    {1\over 2} ( x + |x -y|) \leq x + y \leq 2 ( x + |x -y|).
\end{align*}
In turn we can estimate 
\begin{align*}
\sin \left( {x + y\over 2} \right)  
 \leq x + |x -y| 
\leq {\pi\over 2} \left( \sin\left( x\right) + |x - y|  \right). 
\end{align*}
Likewise, 
\begin{align*}
\sin \left( {x + y\over 2} \right)  
 \geq {2 \over \pi} \left( {x + y \over 2} \right) 
\geq {1\over 2 \pi} \left( x + |x -y| \right)
\geq {1\over 2\pi} \left( \sin\left( x \right) + |x - y|  \right). 
\end{align*}
A similar argument holds in the region ${3 \pi \over4 } \leq x,y \leq \pi$. In particular, noting $\sin({x+y\over2}) = \sin({x-\pi+y-\pi\over2}+\pi) = \sin({(\pi - x )+ (\pi-y)\over2})$, and we can repeat the argument from above by substituting $x$ with $\pi - x$, etc.

\textbf{Case 2.} We now consider the domain $0 \leq x \leq {\pi \over 8}$ and ${\pi \over 4} \leq y \leq \pi$.  In this set 
${\pi \over 8} \leq {x+y\over 2} \leq {9 \pi \over 16}$,  ${\pi \over 8} \leq |x -y| \leq \pi$, and $0\leq \sin(x) \leq {\pi \over 8}$.  In particular,
\[
{1\over 4} = {2 \over \pi} {\pi \over 8} \leq {2 \over \pi}\left( {x + y\over 2} \right) \leq \sin \left( {x + y\over 2} \right)
\leq 1 ;
\]
consequently, 
\begin{align*}
\sin\left( {x + y\over 2} \right)
\leq 1 \leq {8 \over \pi} |x - y| \leq {8 \over \pi} \left( \sin\left( x \right) + |x - y|  \right)
\end{align*}
Likewise,
\begin{align*}
\sin\left( {x + y\over 2} \right)
\geq {1\over 8} + {1\over 8} \geq
{1\over \pi} \sin\left( x\right) + {1\over 8 \pi} |x-y| \geq {1\over 8 \pi}\left( \sin\left( x\right) + |x - y| \right).
\end{align*}
Similar estimates hold in the domain ${7\pi \over 8} \leq x \leq {\pi }$ and ${0} \leq y \leq {3\pi \over 4}$.  

\textbf{Case 3.} We now consider the domain $0 \leq y \leq {\pi \over 8}$ and ${\pi \over 4} \leq x \leq \pi$.  In this set 
${\pi \over 8} \leq {x+y\over 2} \leq {9 \pi \over 16}$ and ${\pi \over 8} \leq |x -y| \leq \pi$. Similar two sided bounds hold.
  
\textbf{Case 4.} Finally, in the set ${\pi \over 8} \leq x, y \leq {7 \pi \over 8}$, we have ${\pi \over 8} \leq {x+ y\over 2} \leq {7 \pi \over 8}$, $0 \leq |x-y| \leq {3 \pi \over 4}$, and ${1\over 4} \leq \sin\left( {x}\right) \leq 1$.  Then
\begin{align*}
\sin\left( {x + y\over 2} \right) \leq {x + y\over 2} \leq {7\pi \over 8} \leq {7\pi \over 2} \sin\left( {x}\right) \leq {7\pi \over 2} \left(\sin\left( {x}\right) + |x - y| \right) 
\end{align*}
and
\begin{align*}
\sin\left( {x + y\over 2} \right) & \geq {2 \over \pi} \left( {x + y\over 2} \right) \geq  
{1\over 8} + {1\over 8} \geq 
{1\over 8} \sin\left( {x}\right) 
+ {1 \over 6\pi} |x-y| \\
& 
\geq {1\over 6\pi}\left(\sin\left( {x }\right) + |x - y| \right) 
\end{align*}

Combining the estimates from Cases 1--4 together yields the estimate for $C = 8\pi$.

\end{proof}

\end{document}
